%========================================================================
\section{Un rôle Product pour compléter la QA}
% TIMING ~4 min  |  La bascule decideur. Une slide = une methode, jamais une liste de taches.
%========================================================================

\begin{frame}{Prendre la main sur la rédaction des tickets}
  \footnotesize Des "QA Checks" à la rédaction complète des tickets $\rightarrow$ rendre les critères d'acceptance explicites.
  \vskip 8pt
  \begin{columns}[T]
    \column{.5\textwidth}
    \begin{block}{Une montée en compétences progressive}
      \footnotesize
      \begin{itemize}
        \item Des tickets simples\dots
        \item \dots aux grosses features découpées équitablement
        \item Jusqu'à animer un \textbf{backlog refinement} complet
      \end{itemize}
    \end{block}
    \column{.47\textwidth}
    \begin{exampleblock}{D'où viennent les besoins ?}
      \footnotesize
      \begin{itemize}
        \item Directement de la Product Owner Polène
        \item Bugs remontés
        \item Besoins d'autres équipes Polène
      \end{itemize}
    \end{exampleblock}
  \end{columns}
  \vskip 16pt
  {\footnotesize $\Rightarrow$~L'objectif est de limiter les allers-retours dev / PO / QA.}
\end{frame}
